%%%%%%%%%%%%%%%%%%%%%%%%%%%%%%%%%%%%%%%%%%%%%%%
%
% Template per Elaborato di Laurea
% DISI - Dipartimento di Ingegneria e Scienza dell’Informazione
%
% update 2015-09-10
%
% Per la generazione corretta del 
% pdflatex nome_file.tex
% bibtex nome_file.aux
% pdflatex nome_file.tex
% pdflatex nome_file.tex
%
%%%%%%%%%%%%%%%%%%%%%%%%%%%%%%%%%%%%%%%%%%%%%%%

% formato FRONTE RETRO
\documentclass[epsfig,a4paper,11pt,titlepage,twoside,openany]{book}
\usepackage{epsfig}
\usepackage{plain}
\usepackage{setspace}
\usepackage[paperheight=29.7cm,paperwidth=21cm,outer=1.5cm,inner=2.5cm,top=2cm,bottom=2cm]{geometry} % per definizione layout
\usepackage{titlesec} % per formato custom dei titoli dei capitoli

%%%%%%%%%%%%%%
% supporto lettere accentate
%
%\usepackage[latin1]{inputenc} % per Windows;
\usepackage[utf8x]{inputenc} % per Linux (richiede il pacchetto unicode);
%\usepackage[applemac]{inputenc} % per Mac.

\singlespacing

\usepackage[italian]{babel}

\begin{document}

  % nessuna numerazione
  \pagenumbering{gobble} 
  \pagestyle{plain}

\thispagestyle{empty}

\begin{center}
  \begin{figure}[h!]
    \centerline{\psfig{file=marchio_unitrento_colore_it_202002.eps,width=0.6\textwidth}}
  \end{figure}

  \vspace{2 cm} 

  \LARGE{Dipartimento di Ingegneria e Scienza dell’Informazione \\ Departement of Information Engineering and Computer Science \\}

  \vspace{1 cm} 
  \Large{Bachelor's Degree in \\ Computer Science}

  \vspace{2 cm} 
  \Large\textsc{Final dissertation\\} 
  \vspace{1 cm} 
  \Huge\textsc{Declarative Specification of Pentesting Strategies for Browser-based Security Protocols:
  the Case Studies of SAML and OAuth/OIDC\\}
  %\Large{\it{Implementation of an automatic testing language and tool to test HTTP browser based protocols}}


  \vspace{1 cm} 
  \begin{tabular*}{\textwidth}{ c @{\extracolsep{\fill}} c }
  \Large{Supervisor} & \Large{Student}\\
  \Large{Silvio Ranise}& \Large{Matteo Bitussi}\\
  \end{tabular*}

  \vspace{0.5 cm} 
  \begin{tabular*}{\textwidth}{ c @{\extracolsep{\fill}} c }
  \Large{Co-Supervisors} & \Large{}\\
  \Large{Andrea Bisegna}& \Large{}\\
  \Large{Roberto Carbone}& \Large{}\\
  \end{tabular*}

  \vspace{2 cm} 

  \Large{Academic year 2021/2022}
  
\end{center}



  \clearpage
 
%%%%%%%%%%%%%%%%%%%%%%%%%%%%%%%%%%%%%%%%%%%%%%%%%%%%%%%%%%%%%%%%%%%%%%%%%%
%%%%%%%%%%%%%%%%%%%%%%%%%%%%%%%%%%%%%%%%%%%%%%%%%%%%%%%%%%%%%%%%%%%%%%%%%%
%% Nota
%%%%%%%%%%%%%%%%%%%%%%%%%%%%%%%%%%%%%%%%%%%%%%%%%%%%%%%%%%%%%%%%%%%%%%%%%%
%% Sezione Ringraziamenti opzionale
%%%%%%%%%%%%%%%%%%%%%%%%%%%%%%%%%%%%%%%%%%%%%%%%%%%%%%%%%%%%%%%%%%%%%%%%%%
%%%%%%%%%%%%%%%%%%%%%%%%%%%%%%%%%%%%%%%%%%%%%%%%%%%%%%%%%%%%%%%%%%%%%%%%%%
  \thispagestyle{empty}

\begin{center}
  {\bf \Huge Ringraziamenti}
\end{center}

\vspace{4cm}


\emph{
  ...thanks to... TODO (in italiano)
}

  \clearpage
  \pagestyle{plain} % nessuna intestazione e pie pagina con numero al centro

  
  % inizio numerazione pagine in numeri arabi
  \mainmatter

%%%%%%%%%%%%%%%%%%%%%%%%%%%%%%%%%%%%%%%%%%%%%%%%%%%%%%%%%%%%%%%%%%%%%%%%%%
%%%%%%%%%%%%%%%%%%%%%%%%%%%%%%%%%%%%%%%%%%%%%%%%%%%%%%%%%%%%%%%%%%%%%%%%%%
%% Nota
%%%%%%%%%%%%%%%%%%%%%%%%%%%%%%%%%%%%%%%%%%%%%%%%%%%%%%%%%%%%%%%%%%%%%%%%%%
%% Si ricorda che il numero massimo di facciate e' 30.
%% Nel conteggio delle facciate sono incluse 
%%   indice
%%   sommario
%%   capitoli
%% Dal conteggio delle facciate sono escluse
%%   frontespizio
%%   ringraziamenti
%%   allegati    
%%%%%%%%%%%%%%%%%%%%%%%%%%%%%%%%%%%%%%%%%%%%%%%%%%%%%%%%%%%%%%%%%%%%%%%%%%
%%%%%%%%%%%%%%%%%%%%%%%%%%%%%%%%%%%%%%%%%%%%%%%%%%%%%%%%%%%%%%%%%%%%%%%%%%

    % indice
    \tableofcontents
    \clearpage
    
    
          
    % gruppo per definizone di successione capitoli senza interruzione di pagina
    \begingroup
      % nessuna interruzione di pagina tra capitoli
      % ridefinizione dei comandi di clear page
      \renewcommand{\cleardoublepage}{} 
      \renewcommand{\clearpage}{} 
      % redefinizione del formato del titolo del capitolo
      % da formato
      %   Capitolo X
      %   Titolo capitolo
      % a formato
      %   X   Titolo capitolo
      
      \titleformat{\chapter}
        {\normalfont\Huge\bfseries}{\thechapter}{1em}{}
        
      \titlespacing*{\chapter}{0pt}{0.59in}{0.02in}
      \titlespacing*{\section}{0pt}{0.20in}{0.02in}
      \titlespacing*{\subsection}{0pt}{0.10in}{0.02in}

      % abstract
      \chapter*{Abstract} % senza numerazione

\addcontentsline{toc}{chapter}{Abstract} % da aggiungere comunque all'indice
% L'abstract deve essere un riassunto dell'introduzione, contesto, motivazione, contributi
This thesis deals with the work started in my internship at Fondazione Bruno Kessler (FBK) in the context of browser-based protocols (based on HTTP) security testing.
In the last years, there has been a constantly rising number of services transitioning from physical to virtual, such as health care and government services, also due to the Covid-19 pandemics. This was done to simplify some procedures and make them faster to be accomplished. However, some security and privacy concerns arise from this transition. A high quantity of sensitive data could be at risk of unwanted access. Software and security engineers have to implement the service's software and manage the data that the software is working with. They will rely on browser-based security protocols to ensure that security requirements such as user identity, the connection between the service and the user, and the integrity of data are met. 

To ensure that all these security requirements are respected during the implementation of the service, there is the need to test the implemented services against the well-known vulnerabilities of the protocols used. Usually, to do that, a software engineer or a security tester, which will be called tester below, will act like an attacker trying to see, edit, and deny access to the data in the services. The tester will use a software to intercept the HTTP messages (from now on messages) between the service and the tester's browser. It will be possible for the tester to edit or remove the messages, simulate an attack, and see if the service is vulnerable. The tester will have to gather all the well-known vulnerabilities related to the used protocols and test them manually. The pentesting is time-consuming and could be done improperly if the tester is not qualified to do it. A software is needed to automatically test the services against the well-known vulnerabilities, which could give a result of whether the services are vulnerable.

The work in this thesis is about the design and implementation of the automatic HTTP pentesting software (from now on tool) and its relative declarative specification language. The tool addresses the challenges highlighted above, such as the automation of the pentesting phase to improve its precision and the design of a language for the specification of the tests. Other software has been evaluated before the start of this work, some of them missed the automation part, others were limited in terms of the possible security tests (from now on tests) to be done. For these reasons, a tool and language have been designed and implemented. First, a declarative specification of pentesting strategies has been designed, simultaneously with the design of an approach to execute the pentesting strategies. Next, the implementation of a tool to execute the pentesting strategies has been done. To test the correctness of the implemented tool and the expressiveness of the proposed specification language, a test definition for \gls{OAuth}/OIDC, a web security protocol, has been written in the implemented language and used during development.

The implemented tool and language have proved to be working in two different use cases, with OAuth/OIDC and SAML. The results obtained by the tool have also been confirmed by comparison with the results of similar software; they have proved to be correct. Lastly, the objectives fixed at the beginning of the work were accomplished, the software and language created are valuable tools to help the tester discover vulnerabilities in services. The definition of an OAuth/OIDC and SAML test suite covering the well-known vulnerabilities is useful as the tester does not need to write them every time to test them. Depending on the service to test, what has to be changed are the automation actions to be executed during testing, as the web pages may be different, so the actions to be accomplished have to be different. Future work has been planned to extend the tool with other functionalities and support other protocols. 








      % background
      \chapter{Background}
\label{Background}

The idea of a pentesting tool used to test SSO implementations such as OAuth, OIDC and SAML was previously discussed by my colleagues Stefano Facchini \cite{stefano_faccini}, Claudio Grisenti \cite{claudio_grisenti} and Wendy Barreto \cite{wendy_barreto}, which developed a plugin in Burp, that is a toolkit for web security testing.

\section{The old plugin}
The plugin work by executing a browser session and by doing some automatic actions defined in a track. All the messages received and sent by the browser are intercepted by the Burp's proxy. The plugin checks the messages and based on the test defined in the plugin tells if there is a vulnerability or not.
Starting from the first version of Stefano, to the last of Wendy, the plugin was improved. In the first two versions of Stefano and Claudio the plugin was static, in a way that ony the supported tests could be executed, with little settings to change. This version of the plugin worked well, but as said the tests could not be customized.

\section{Last plugin version}
Wendy improved the plugin by removing the staticity of the test, adding the possibility to customly define all of the tests by some specific JSON objects. \\
This is the last version of the plugin which i started working to. The plugin supported the definition of static and dynamic tests: static tests are tests where the messages are not edited, dynamic tests are tests where there could be an edit of one or more messages. The tests where working, but were very limited in the types of actions that could be defined. For example:
\begin{itemize}
    \item Limited oracle for the verification of active tests
    \item The filtering of which message to check or edit for static tests was limited: ( only "Authorization grant message", "Response messages", "Request messages" and "All messages")
    \item Only regex were supported to search something in a message
    \item Unable to work over encoded parameters
    \item Impossibility of doing multiple operations on a single message (verify)
    \item Impossibility of saving a parameter and using it somewhere else
\end{itemize}
Moreover, the language was thought with OIDC and OAuth in mind, other SSO protocols such as SAML could not be tested, because of the fact that SAML parameters are encoded, so editing them or verifying them is not possible.




      % sommario
      \chapter*{Introduction} % senza numerazione
\label{Introduction}

\addcontentsline{toc}{chapter}{Sommario} % da aggiungere comunque all'indice

Sommario è un breve riassunto del lavoro svolto dove si descrive l'obiettivo, l'oggetto della tesi, le 
metodologie e le tecniche usate, i dati elaborati e la spiegazione delle conclusioni alle quali siete arrivati.  

Il sommario dell’elaborato consiste al massimo di 3 pagine e deve contenere le seguenti informazioni:
\begin{itemize}
  \item contesto e motivazioni 
  \item breve riassunto del problema affrontato
  \item tecniche utilizzate e/o sviluppate
  \item risultati raggiunti, sottolineando il contributo personale del laureando/a
\end{itemize}


%%%%%%%%%%%%%%%%%%%%%%%%%%%%%%%%%%%%%%%%%%%%%%%%%%%%%%%%%%%%%%%%%%%%%%%%%%
%%%%%%%%%%%%%%%%%%%%%%%%%%%%%%%%%%%%%%%%%%%%%%%%%%%%%%%%%%%%%%%%%%%%%%%%%%
%% Nota
%%%%%%%%%%%%%%%%%%%%%%%%%%%%%%%%%%%%%%%%%%%%%%%%%%%%%%%%%%%%%%%%%%%%%%%%%%
%% Sommario e' un breve riassunto del lavoro svolto dove si descrive 
%% l’obiettivo, l’oggetto della tesi, le metodologie e 
%% le tecniche usate, i dati elaborati e la spiegazione delle conclusioni 
%% alle quali siete arrivati.
%% Il sommario dell’elaborato consiste al massimo di 3 pagine e deve contenere le seguenti informazioni: 
%%   contesto e motivazioni
%%   breve riassunto del problema affrontato
%%   tecniche utilizzate e/o sviluppate
%%   risultati raggiunti, sottolineando il contributo personale del laureando/a
%%%%%%%%%%%%%%%%%%%%%%%%%%%%%%%%%%%%%%%%%%%%%%%%%%%%%%%%%%%%%%%%%%%%%%%%%%
%%%%%%%%%%%%%%%%%%%%%%%%%%%%%%%%%%%%%%%%%%%%%%%%%%%%%%%%%%%%%%%%%%%%%%%%%%      
      
      %%%%%%%%%%%%%%%%%%%%%%%%%%%%%%%%
      % lista dei capitoli
      %
      % \input oppure \include
      %
      \chapter{Changes to the language and plugin}
\label{cha:intro}

\section{A more abstract Language}
\label{sec:context}

I decided to start from the last version of the language and to completely rethink it, whith a more abstract approach, in a way that almost any type of protocol that uses HTTP could be tested. I started gathering all types of tests that could be wanted to be executed, searching for all the possible actions that could be done to a message.

\begin{center}
    \begin{tabular}{ |c|c|c| }
        \hline
        Action                                     & Old language         & New language                             \\
        \hline
        Custom message filtering                   & Only on active tests & supported                                \\
        Edit string                                & only by regex        & supported with regex and check construct \\
        Remove string                              & only by regex        & supported with regex and check construct \\
        Add string                                 & not supported        & supported                                \\
        Check parameter                            & only with regex      & with regex and check construct           \\
        Multiple checks in single message          & not supported        & supported                                \\
        Saving and reusing of strings and messages & not supported        & supported                                \\
        Multiple sessions in single test           & not supported        & supported                                \\
        Custom oracle definition                   & not supported        & by using regex and checks                \\

        \hline
    \end{tabular}
\end{center}
      \chapter{Proin rhoncus a sapien in.}
\label{cha:789}
Lorem ipsum dolor sit amet, consectetur adipiscing elit. Donec sed nunc orci. Aliquam nec nisl vitae sapien pulvinar dictum quis non urna. Suspendisse at dui a erat aliquam vestibulum. Quisque ultrices pellentesque pellentesque. Pellentesque egestas quam sed blandit tempus. Sed congue nec risus posuere euismod. Maecenas ut lacus id mauris sagittis egestas a eu dui. Class aptent taciti sociosqu ad litora torquent per conubia nostra, per inceptos himenaeos. Pellentesque at ultrices tellus. Ut eu purus eget sem iaculis ultricies sed non lorem. Curabitur gravida dui eget ex vestibulum venenatis. Phasellus gravida tellus velit, non eleifend justo lobortis eget. 


\section{Cras in aliquam quam, et}
\label{sec:456}
Lorem ipsum dolor sit amet, consectetur adipiscing elit. Donec sed nunc orci. Aliquam nec nisl vitae sapien pulvinar dictum quis non urna. Suspendisse at dui a erat aliquam vestibulum. Quisque ultrices pellentesque pellentesque. Pellentesque egestas quam sed blandit tempus. Sed congue nec risus posuere euismod. Maecenas ut lacus id mauris sagittis egestas a eu dui. Class aptent taciti sociosqu ad litora torquent per conubia nostra, per inceptos himenaeos. Pellentesque at ultrices tellus. Ut eu purus eget sem iaculis ultricies sed non lorem. Curabitur gravida dui eget ex vestibulum venenatis. Phasellus gravida tellus velit, non eleifend justo lobortis eget.


\subsection{Sed pulvinar placerat enim, a}
\label{sec:00456}
Lorem ipsum dolor sit amet, consectetur adipiscing elit. Donec sed nunc orci. Aliquam nec nisl vitae sapien pulvinar dictum quis non urna. Suspendisse at dui a erat aliquam vestibulum. Quisque ultrices pellentesque pellentesque. Pellentesque egestas quam sed blandit tempus. Sed congue nec risus posuere euismod. Maecenas ut lacus id mauris sagittis egestas a eu dui. Class aptent taciti sociosqu ad litora torquent per conubia nostra, per inceptos himenaeos. Pellentesque at ultrices tellus. Ut eu purus eget sem iaculis ultricies sed non lorem. Curabitur gravida dui eget ex vestibulum venenatis. Phasellus gravida tellus velit, non eleifend justo lobortis eget.


\section{Vivamus hendrerit imperdiet ex. Vivamus}
\label{sec:123}
Lorem ipsum dolor sit amet, consectetur adipiscing elit. Donec sed nunc orci. Aliquam nec nisl vitae sapien pulvinar dictum quis non urna. Suspendisse at dui a erat aliquam vestibulum. Quisque ultrices pellentesque pellentesque. Pellentesque egestas quam sed blandit tempus. Sed congue nec risus posuere euismod. Maecenas ut lacus id mauris sagittis egestas a eu dui. Class aptent taciti sociosqu ad litora torquent per conubia nostra, per inceptos himenaeos. Pellentesque at ultrices tellus. Ut eu purus eget sem iaculis ultricies sed non lorem. Curabitur gravida dui eget ex vestibulum venenatis. Phasellus gravida tellus velit, non eleifend justo lobortis eget.



      \chapter{Conclusions}
\label{cha:Conclusions}
Lorem ipsum dolor sit amet, consectetur adipiscing elit. Donec sed nunc orci. Aliquam nec nisl vitae sapien pulvinar dictum quis non urna. Suspendisse at dui a erat aliquam vestibulum. Quisque ultrices pellentesque pellentesque. Pellentesque egestas quam sed blandit tempus. Sed congue nec risus posuere euismod. Maecenas ut lacus id mauris sagittis egestas a eu dui. Class aptent taciti sociosqu ad litora torquent per conubia nostra, per inceptos himenaeos. Pellentesque at ultrices tellus. Ut eu purus eget sem iaculis ultricies sed non lorem. Curabitur gravida dui eget ex vestibulum venenatis. Phasellus gravida tellus velit, non eleifend justo lobortis eget.



      %\input{capitolo4}
      
      
    \endgroup


    % bibliografia in formato bibtex
    %
    % aggiunta del capitolo nell'indice
    \addcontentsline{toc}{chapter}{Bibliografia}
    % stile con ordinamento alfabetico in funzione degli autori
    \bibliographystyle{plain}
    \bibliography{biblio}
%%%%%%%%%%%%%%%%%%%%%%%%%%%%%%%%%%%%%%%%%%%%%%%%%%%%%%%%%%%%%%%%%%%%%%%%%%
%%%%%%%%%%%%%%%%%%%%%%%%%%%%%%%%%%%%%%%%%%%%%%%%%%%%%%%%%%%%%%%%%%%%%%%%%%
%% Nota
%%%%%%%%%%%%%%%%%%%%%%%%%%%%%%%%%%%%%%%%%%%%%%%%%%%%%%%%%%%%%%%%%%%%%%%%%%
%% Nella bibliografia devono essere riportati tutte le fonti consultate 
%% per lo svolgimento della tesi. La bibliografia deve essere redatta 
%% in ordine alfabetico sul cognome del primo autore. 
%% 
%% La forma della citazione bibliografica va inserita secondo la fonte utilizzata:
%% 
%% LIBRI
%% Cognome e iniziale del nome autore/autori, la data di edizione, titolo, casa editrice, eventuale numero dell’edizione. 
%% 
%% ARTICOLI DI RIVISTA
%% Cognome e iniziale del nome autore/autori, titolo articolo, titolo rivista, volume, numero, numero di pagine.
%% 
%% ARTICOLI DI CONFERENZA
%% Cognome e iniziale del nome autore/autori (anno), titolo articolo, titolo conferenza, luogo della conferenza (città e paese), date della conferenza, numero di pagine. 
%% 
%% SITOGRAFIA
%% La sitografia contiene un elenco di indirizzi Web consultati e disposti in ordine alfabetico. 
%% E’ necessario:
%%   Copiare la URL (l’indirizzo web) specifica della pagina consultata
%%   Se disponibile, indicare il cognome e nome dell’autore, il titolo ed eventuale sottotitolo del testo
%%   Se disponibile, inserire la data di ultima consultazione della risorsa (gg/mm/aaaa).    
%%%%%%%%%%%%%%%%%%%%%%%%%%%%%%%%%%%%%%%%%%%%%%%%%%%%%%%%%%%%%%%%%%%%%%%%%%
%%%%%%%%%%%%%%%%%%%%%%%%%%%%%%%%%%%%%%%%%%%%%%%%%%%%%%%%%%%%%%%%%%%%%%%%%%

    \titleformat{\chapter}
        {\normalfont\Huge\bfseries}{Allegato \thechapter}{1em}{}
    % sezione Allegati - opzionale
    \appendix
    \chapter{PKCE test Example complete}
\label{chap:PKCE_complete}
\begin{lstlisting}[language=json]
{
    "test suite": {
        "name": "OAuth Active tests",
        "description": "A series of tests to test OAuth's well-known vulnerabilities",
        "filter messages": true
    },
    "tests": [     
        {
            "test": {
                "name": "PKCE Downgrade",
                "description": "Tries to remove code_challenge parameter",
                "type": "active",
                "sessions": [
                    "s1"
                ],
                "operations": [
                    {
                        "session": "s1",
                        "action": "start"
                    },
                    {
                        "action": "intercept",
                        "from session": "s1",
                        "then": "forward",
                        "message type": "authorization request",
                        "preconditions": [
                            {
                                "in": "url",
                                "check param": "code_challenge",
                                "is present": true
                            }
                        ],
                        "message operations": [
                            {
                                "from": "url",
                                "remove parameter": "code_challenge"
                            }
                        ]
                    }
                ],
                "result": "incorrect flow s1"
            }
        }
    ]
}
\end{lstlisting}

\chapter{Language comparison}
\label{attachment:languages_comparison}

\begin{center}
    \begin{tabular}{ |c|c|c| }
        \hline
        Action                                    & Old language         & New language                             \\
        \hline\hline
        Custom message filtering                  & Only on active tests & supported                                \\
        Edit string                               & only by regex        & supported with regex and check construct \\
        Remove string                             & only by regex        & supported with regex and check construct \\
        Add string                                & not supported        & supported                                \\
        Check parameter                           & only with regex      & with regex and check construct           \\
        Multiple operations in single message     & not supported        & supported                                \\
        Saving and reusing of values and messages & not supported        & supported                                \\
        Multiple sessions in single test          & not supported        & supported                                \\
        Custom oracle definition                  & not supported        & by using regex and checks                \\
        \hline
    \end{tabular}
\end{center}






\end{document}
